\chapter{Gradient flow of the empirical risk: function-space dynamics and lazy training}\label{chap:ntkgf}

This chapter records the model-independent theory of gradient flow near a random initialization, following
Jacot et al.\ (2018) and Lee et al.\ (2019). Everything is stated for an arbitrary parametrized model
$f(x;\theta)$ with $\theta\in\mathbb R^P$, trained on a finite dataset by gradient flow on an empirical risk.
The shallow and two-layer networks of the other chapters enter only through hypotheses on the output
Jacobian (a bound $M$ and a Lipschitz constant $L_J$) and on the initial kernel (a spectral gap).
The Lean files are in \texttt{Optimization/NTK/Shallow/DatasetNTK.lean} and
\texttt{Optimization/NTK/Training/GradientFlow/} (modules \texttt{Dynamics}, \texttt{Convergence},
\texttt{LinearDynamics}, \texttt{KernelStability}, \texttt{Bootstrap}, \texttt{InfiniteWidth},
\texttt{ODEStability}, \texttt{AffineDynamics}, \texttt{TestPrediction}).

\medskip\noindent\textbf{Notation.}
In this chapter $m$ is the \emph{number of training points} and $P$ the number of parameters (the network
width, written $n$ in the two-layer chapters, enters only through the constants $M,L_J,C$).
The training inputs are $x^1,\dots,x^m$ and the labels form $y\in\mathbb R^m$. Vectors in $\mathbb R^m$ and
$\mathbb R^P$ carry the Euclidean norm, and matrices the Frobenius norm. The ODE tools of
Chapter~\ref{chap:ntkf} (Gr\"onwall, bootstrap, existence of forward flows) are used throughout.


\section{The dataset-level objects}\label{sec:ntkgf:setting}

\begin{definition}[Training outputs, residual, loss]\label{def:ntkgf:outputs}
  \lean{NTK.trainingOutputs, NTK.trainingResidual, NTK.mseLoss}
  \leanok
For a model $f:\mathcal X\times\mathbb R^P\to\mathbb R$ and a dataset $x^1,\dots,x^m$ define
\[
  f(\theta):=\bigl(f(x^\alpha;\theta)\bigr)_{\alpha\le m}\in\mathbb R^m,\qquad
  r(\theta):=f(\theta)-y,\qquad
  L(\theta):=\frac1{2m}\|r(\theta)\|^2=\frac1{2m}\sum_\alpha\bigl(f(x^\alpha;\theta)-y^\alpha\bigr)^2 .
\]
\end{definition}

\begin{definition}[Tangent features, output Jacobian, empirical NTK matrix]\label{def:ntkgf:jacobian}
  \uses{def:ntkgf:outputs}
  \lean{NTK.tangentFeature, NTK.outputJacobian, NTK.empiricalNTKMatrix}
  \leanok
The \emph{tangent feature} of an input $x$ is $\nabla_\theta f(x;\theta)\in\mathbb R^P$. The \emph{output
Jacobian} $J(\theta)\in\mathbb R^{m\times P}$ has $\alpha$-th row $\nabla_\theta f(x^\alpha;\theta)^\top$. The
\emph{empirical NTK matrix} is the Gram matrix of the tangent features,
\[
  K(\theta):=J(\theta)J(\theta)^\top\in\mathbb R^{m\times m},\qquad
  K^{\alpha\beta}(\theta)=\bigl\langle\nabla_\theta f(x^\alpha;\theta),\nabla_\theta f(x^\beta;\theta)\bigr\rangle .
\]
\end{definition}

\begin{lemma}[Basic properties of the empirical NTK matrix]\label{lem:ntkgf:ntkPsd}
  \uses{def:ntkgf:jacobian}
  \lean{NTK.empiricalNTKMatrix_apply, NTK.empiricalNTKMatrix_posSemidef,
        NTK.empiricalNTKMatrix_quad_form, NTK.empiricalNTKMatrix_quad_form_nonneg}
  \leanok
\begin{enumerate}
  \item $K^{\alpha\beta}(\theta)=\langle\nabla f(x^\alpha;\theta),\nabla f(x^\beta;\theta)\rangle$.
  \item $K(\theta)$ is positive semidefinite, and $v^\top K(\theta)v=\|J(\theta)^\top v\|^2\ge0$ for every $v$.
\end{enumerate}
(The row of this matrix belonging to a \emph{test} input is the cross-kernel of Definition~\ref{def:ntkgf:crossKernel} below.)
\end{lemma}

\begin{proof}
  \leanok
(1) is the definition of the product $JJ^\top$ in terms of the rows. For (2), $v^\top JJ^\top v=\|J^\top v\|^2$.
\end{proof}

\subsection{The train--test cross-kernel}\label{sec:ntkgf:crossKernel}

The empirical NTK matrix $K(\theta)$ only compares \emph{training} inputs with each other. To say how a gradient step on the training loss changes the network's output at an input $x$ that is \emph{not} in the training set (a \emph{test} input), one needs the kernel between $x$ and each training point.
That object is the cross-kernel.

\begin{definition}[Train--test cross-kernel]\label{def:ntkgf:crossKernel}
  \uses{def:ntkgf:jacobian}
Fix a test input $x$ and write $g_\theta(x):=\nabla_\theta f(x;\theta)\in\mathbb R^P$ for its tangent feature (Definition~\ref{def:ntkgf:jacobian}). The \emph{cross-kernel vector} of $x$ at the parameter $\theta$ is the vector $k_\theta(x)\in\mathbb R^m$ with entries
\[
  k_\theta(x)_\alpha\;:=\;\bigl\langle\nabla_\theta f(x;\theta),\,\nabla_\theta f(x^\alpha;\theta)\bigr\rangle ,\qquad\alpha=1,\dots,m,
\]
that is, the inner product of the test point's tangent feature with each training point's tangent feature. It is the NTK evaluated on the pairs $(x,x^\alpha)$.
Collecting it with the \emph{test norm} $q_\theta(x):=\|g_\theta(x)\|^2=\langle g_\theta(x),g_\theta(x)\rangle$, the empirical NTK matrix of the \emph{extended dataset} $(x^1,\dots,x^m,x)$ has the block form
\[
  K^{\mathrm{ext}}(\theta)=\begin{pmatrix}K(\theta)&k_\theta(x)\\ k_\theta(x)^\top&q_\theta(x)\end{pmatrix}.
\]
In Lean there is no separate definition: the cross-kernel is always written out as $\langle\texttt{tangentFeature}\,f\,x\,\theta,\texttt{tangentFeature}\,f\,(X\,\alpha)\,\theta\rangle$, or as the vector $J(\theta)\,\nabla f(x;\theta)$, and the two lemmas below identify these descriptions.
\end{definition}

\begin{lemma}[Three descriptions of the cross-kernel]\label{lem:ntkgf:crossKernelFormulas}
  \uses{def:ntkgf:crossKernel, lem:ntkgf:ntkPsd}
  \lean{NTK.empiricalNTKMatrix_snoc_last, NTK.outputJacobian_mulVec_tangentFeature}
  \leanok
For a test input $x$, every $\theta$ and every training index $\alpha$:
\begin{enumerate}
  \item (last row of an extended NTK) In the empirical NTK matrix of the extended dataset $(x^1,\dots,x^m,x)$, the entry in the last row and column $\alpha$ is $\langle\nabla f(x;\theta),\nabla f(x^\alpha;\theta)\rangle=k_\theta(x)_\alpha$, and the last diagonal entry is $\|\nabla f(x;\theta)\|^2=q_\theta(x)$.
  \item (matrix--vector form) $k_\theta(x)=J(\theta)\,\nabla_\theta f(x;\theta)$, where $J(\theta)$ is the output Jacobian of the \emph{training} dataset.
\end{enumerate}
\end{lemma}

\begin{proof}
  \uses{lem:ntkgf:ntkPsd}
  \leanok
(1) is Lemma~\ref{lem:ntkgf:ntkPsd}(1) applied to the extended dataset, whose last point is $x$ (and the first $m$ points are the training points). (2) The $\alpha$-th entry of $J(\theta)g$ is the inner product of the $\alpha$-th row of $J(\theta)$, namely $\nabla f(x^\alpha;\theta)^\top$, with $g=\nabla f(x;\theta)$.
\end{proof}

\begin{remark}[What the cross-kernel measures]\label{rmk:ntkgf:crossKernelMeaning}
  \uses{def:ntkgf:crossKernel, thm:ntkgf:outputODE, thm:ntkgf:mseGradient}
Gradient flow moves the parameter along $\theta'=-\nabla L(\theta)=-\frac1m\sum_\beta r^\beta(\theta)\,\nabla f(x^\beta;\theta)$. The output at the test input therefore changes at the rate
\[
  \partial_t f(x;\theta(t))=\bigl\langle\nabla f(x;\theta),\theta'\bigr\rangle=-\frac1m\sum_{\beta=1}^m\bigl\langle\nabla f(x;\theta),\nabla f(x^\beta;\theta)\bigr\rangle\,r^\beta=-\frac1m\,k_\theta(x)^\top r(t),
\]
which is the exact analogue, for a test point, of the training-output equation $\partial_tf(t)=-\frac1mK_tr(t)$ of Theorem~\ref{thm:ntkgf:outputODE}. Entry $\alpha$ of $k_\theta(x)$ says how strongly the residual on training point $\alpha$ pushes the output at $x$; if $x=x^\alpha$ is itself a training input the vector becomes row $\alpha$ of $K(\theta)$.
This is the same chain-rule computation as Steps~1--3 of Theorem~\ref{thm:ntkgf:outputODE}; the Lean development states Step~1 for an arbitrary dataset (so a one-point dataset $\{x\}$ gives the test-output derivative, \texttt{NTK.hasDerivAt\_trainingOutputs\_coord}), and works with the \emph{frozen} feature $g=\nabla f(x;\theta_0)$ in
Lemma~\ref{lem:ntkgf:predictionDerivative} below rather than stating the test-point equation separately.
\end{remark}

\begin{theorem}[Gradient of the MSE loss]\label{thm:ntkgf:mseGradient}
  \uses{def:ntkgf:outputs, def:ntkgf:jacobian}
  \lean{NTK.hasFDerivAt_generalizedLoss_term, NTK.hasFDerivAt_sq_diff, NTK.mseLoss_eq_sum,
        NTK.hasGradientAt_mseLoss, NTK.gradient_mseLoss, NTK.gradient_mseLoss_apply_j,
        NTK.gradient_mseLoss_eq_mulVec, NTK.gdIterate_mseLoss_succ, NTK.gradient_mseLoss_norm_le}
  \leanok
If every $\theta'\mapsto f(x^\alpha;\theta')$ is differentiable at $\theta$, then
\[
  \nabla_\theta L(\theta)=\frac1m\sum_{\alpha=1}^mr^\alpha(\theta)\,\nabla_\theta f(x^\alpha;\theta)=\frac1mJ(\theta)^\top r(\theta),
\]
coordinatewise $[\nabla L(\theta)]_j=\frac1m[J(\theta)^\top r(\theta)]_j$. The gradient-descent iteration with step size $\eta$
reads $\theta_{k+1}=\theta_k-\eta\nabla L(\theta_k)$, and the speed is bounded by
\[
  \|\nabla_\theta L(\theta)\|\le\frac1m\,\|J(\theta)\|_{\mathrm F}\,\|r(\theta)\| .
\]
\end{theorem}

\begin{proof}
  \leanok
By the chain rule, the map $\theta\mapsto\ell(g(\theta))$ has gradient $\ell'(g(\theta))\nabla g(\theta)$; for
$\ell(u)=(u-c)^2$ this is $2(g(\theta)-c)\nabla g(\theta)$. Summing over $\alpha$ with the weight $1/(2m)$ gives
the formula. The speed bound is Lemma~\ref{lem:ntkf:mulVec}: $\|J^\top r\|\le\|J\|_{\mathrm F}\|r\|$.
\end{proof}

\begin{definition}[Forward gradient-flow trajectory]\label{def:ntkgf:forwardGF}
  \lean{ConvexOpt.ForwardGFTrajectory}
  \leanok
A curve $\theta:\mathbb R\to\mathbb R^P$ is a \emph{forward gradient flow} of a loss $L$ started at $\theta_0$ if
$\theta(0)=\theta_0$, $\theta$ is continuous on $[0,\infty)$, and
$\theta'(t)=-\nabla L(\theta(t))$ for every $t>0$. Only forward time is used: for an unbounded activation a
solution on $[0,\infty)$ need not extend backward.
\end{definition}


\section{Function-space dynamics under gradient flow}\label{sec:ntkgf:dynamics}

\begin{theorem}[Output and residual ODE, Proposition~2.15]\label{thm:ntkgf:outputODE}
  \uses{def:ntkgf:jacobian, thm:ntkgf:mseGradient}
  \lean{NTK.euclideanSpace_inner_eq_sum, NTK.hasDerivAt_trainingOutputs_coord,
        NTK.hasDerivAt_trainingOutputs_coord_sum, NTK.hasDerivAt_euclideanSpace,
        NTK.gradient_flow_output_coord_deriv_eq_inner_grad,
        NTK.gradient_flow_output_coord_deriv_eq_sum_inner, NTK.gradient_flow_output_coord_ode,
        NTK.gradient_flow_output_vector_ode, NTK.gradient_flow_output_vector_ode_sub_y,
        NTK.gradient_flow_residual_vector_ode}
  \leanok
Let $\theta$ be differentiable at $t$ with $\theta'(t)=-\nabla L(\theta(t))$ and suppose each output is differentiable in
$\theta$ at $\theta(t)$. Then the training outputs and residual are differentiable at $t$ with
\[
  \partial_tf(t)=-\frac1mK_t\,r(t)=-\frac1mK_t\bigl(f(t)-y\bigr),\qquad
  \partial_tr(t)=-\frac1mK_t\,r(t),\qquad K_t:=K(\theta(t)),
\]
that is, coordinatewise $\partial_tf^\alpha(t)=-\frac1m\sum_\beta K_t^{\alpha\beta}r^\beta(t)$.
\end{theorem}

\begin{proof}
  \leanok
The proof is a five-step chain-rule computation.

\emph{Step 1.} By the multivariate chain rule,
$\partial_tf^\alpha(t)=\sum_{j}\partial_{\theta_j}f(x^\alpha;\theta(t))\,\partial_t\theta_j(t)=\langle\nabla f^\alpha,\partial_t\theta\rangle$.

\emph{Step 2.} Insert the flow law $\partial_t\theta=-\nabla L(\theta)$.

\emph{Step 3.} Insert $\nabla L=\frac1m\sum_\beta r^\beta\nabla f^\beta$ (Theorem~\ref{thm:ntkgf:mseGradient}), giving
$-\frac1m\sum_\beta\langle\nabla f^\alpha,\nabla f^\beta\rangle r^\beta$.

\emph{Step 4.} Recognise the entries $K_t^{\alpha\beta}$ of the empirical NTK matrix.

\emph{Step 5.} Assemble the coordinates into the vector identity, using that a curve in $\mathbb R^m$ is differentiable iff each coordinate is.
The residual differs from the output by the constant $y$.
\end{proof}

\begin{definition}[Generalized empirical risk]\label{def:ntkgf:generalizedRisk}
  \uses{def:ntkgf:outputs}
  \lean{NTK.generalizedEmpiricalRisk, NTK.generalizedResidual}
  \leanok
For a pointwise loss $\ell:\mathbb R\times\mathbb R\to\mathbb R$ differentiable in its first argument, define
$L_\ell(\theta)=\frac1m\sum_\alpha\ell(f(x^\alpha;\theta),y^\alpha)$ and the generalized residual
$r^\alpha(\theta)=\partial_f\ell(f(x^\alpha;\theta),y^\alpha)$. Examples:
\begin{center}
\begin{tabular}{lll}
  squared loss & $\ell(f,y)=\tfrac12(f-y)^2$ & $r^\alpha=f^\alpha-y^\alpha$,\\
  logistic loss & $\ell(f,y)=\log(1+e^{-yf})$ & $r^\alpha=-y^\alpha\,\operatorname{sigmoid}(-y^\alpha f^\alpha)$,\\
  exponential loss & $\ell(f,y)=e^{-yf}$ & $r^\alpha=-y^\alpha e^{-y^\alpha f^\alpha}$.
\end{tabular}
\end{center}
\end{definition}

\begin{lemma}[Residuals of the standard losses]\label{lem:ntkgf:exampleLosses}
  \uses{def:ntkgf:generalizedRisk}
  \lean{NTK.hasDerivAt_squaredLoss, NTK.hasDerivAt_logisticLoss, NTK.hasDerivAt_exponentialLoss}
  \leanok
For all real $f,y$: $\frac{d}{df}\tfrac12(f-y)^2=f-y$; $\frac{d}{df}\log(1+e^{-yf})=-y\,\operatorname{sigmoid}(-yf)$;
$\frac{d}{df}e^{-yf}=-ye^{-yf}$.
\end{lemma}

\begin{proof}
  \leanok
Direct differentiation; for the logistic loss, $\frac{d}{df}\log(1+e^{-yf})=\frac{-ye^{-yf}}{1+e^{-yf}}=-y\,\operatorname{sigmoid}(-yf)$.
\end{proof}

\begin{theorem}[Output dynamics for a general loss]\label{thm:ntkgf:generalOutputODE}
  \uses{def:ntkgf:generalizedRisk, thm:ntkgf:outputODE}
  \lean{NTK.gradient_generalizedRisk, NTK.gradient_flow_generalizedOutput_coord_deriv_eq_inner_grad,
        NTK.gradient_flow_generalizedOutput_coord_deriv_eq_sum_inner,
        NTK.gradient_flow_generalizedOutput_coord_ode, NTK.gradient_flow_generalizedOutput_vector_ode}
  \leanok
Under the differentiability hypotheses on $f$ and $\ell$, $\nabla L_\ell(\theta)=\frac1m\sum_\alpha r^\alpha(\theta)\nabla f(x^\alpha;\theta)$, and
along a gradient flow of $L_\ell$,
\[
  \partial_tf(t)=-\frac1mK_t\,r(t),\qquad r^\alpha(t)=\partial_f\ell\bigl(f^\alpha(t),y^\alpha\bigr).
\]
Proposition~2.15 (squared loss) is the special case $\ell=\frac12(f-y)^2$.
\end{theorem}

\begin{proof}
  \leanok
The proof of Theorem~\ref{thm:ntkgf:outputODE} goes through verbatim with $r^\alpha=f^\alpha-y^\alpha$ replaced by the
generalized residual; the gradient formula is the chain rule applied to each term $\ell(f(x^\alpha;\theta),y^\alpha)$.
\end{proof}

\begin{theorem}[Risk dissipation identity, Proposition~2.17]\label{thm:ntkgf:riskDissipation}
  \uses{thm:ntkgf:generalOutputODE, lem:ntkgf:ntkPsd}
  \lean{NTK.hasDerivAt_generalizedEmpiricalRisk_coord_sum, NTK.risk_dissipation_identity,
        NTK.risk_dissipation_nonpos, NTK.risk_dissipation_le_of_rayleighRitz}
  \leanok
Along a gradient flow of $L_\ell$,
\[
  \partial_tL_\ell(\theta(t))=\frac1m\,r(t)^\top\partial_tf(t)=-\frac1{m^2}\,r(t)^\top K_t\,r(t)\;\le\;0 .
\]
If moreover $\lambda\|v\|^2\le v^\top K_tv$ for all $v$ (a Rayleigh--Ritz bound), then
$\partial_tL_\ell(\theta(t))\le-\frac{\lambda}{m^2}\|r(t)\|^2$.
\end{theorem}

\begin{proof}
  \leanok
\emph{Step 1.} By the chain rule $\partial_tL_\ell=\frac1m\sum_\alpha r^\alpha\partial_tf^\alpha=\frac1mr^\top\partial_tf$.
\emph{Step 2.} Insert $\partial_tf=-\frac1mK_tr$ (Theorem~\ref{thm:ntkgf:generalOutputODE}). Nonpositivity follows since $K_t\succeq0$
(Lemma~\ref{lem:ntkgf:ntkPsd}); the last bound is the Rayleigh--Ritz assumption applied to $v=r(t)$.
\end{proof}

\begin{theorem}[Kinetic energy of the flow]\label{thm:ntkgf:kineticEnergy}
  \uses{def:ntkgf:forwardGF, thm:ntkgf:riskDissipation}
  \lean{NTK.norm_sq_sum_smul_eq_dotProduct, NTK.norm_sq_gradient_generalizedRisk,
        NTK.norm_sq_gradient_mseLoss, NTK.norm_deriv_sq_eq_quadratic_form_of_forwardGF,
        NTK.mseLoss_sub_eq_integral_quadratic_form, NTK.hasDerivAt_coord_of_forwardGF}
  \leanok
For every $\theta$, $\|\nabla L(\theta)\|^2=\frac1{m^2}r(\theta)^\top K(\theta)r(\theta)$ (for $L_\ell$ and in particular for the MSE loss). Along a
forward gradient flow of the MSE loss, for $t>0$,
\[
  \|\theta'(t)\|^2=\frac1{m^2}\,r(t)^\top K_tr(t),\qquad
  L(\theta(0))-L(\theta(T))=\int_0^T\frac1{m^2}r(t)^\top K_tr(t)\,dt=\int_0^T\|\theta'(t)\|^2\,dt,
\]
and every parameter coordinate moves as $\partial_t\theta_k=-\frac1m[J^\top r]_k=-\frac1m\sum_\alpha r^\alpha\partial_{\theta_k}f^\alpha$.
\end{theorem}

\begin{proof}
  \leanok
The squared norm of a linear combination $\|\sum_\alpha c_\alpha v_\alpha\|^2$ equals $c^\top Gc$, where $G$ is the Gram matrix of the vectors
$v_\alpha$; apply this to $\nabla L=\frac1m\sum r^\alpha\nabla f^\alpha$. Along the flow $\theta'=-\nabla L$, so $\|\theta'\|^2=\|\nabla L\|^2$.
By the chain rule $\partial_tL(\theta(t))=-\|\nabla L(\theta(t))\|^2=-\|\theta'(t)\|^2$ (equivalently, Theorem~\ref{thm:ntkgf:riskDissipation}), and
integrating over $[0,T]$ gives the energy identity. The coordinate equation reads off the Jacobian entries in $\theta'=-\frac1mJ^\top r$.
\end{proof}


\section{Exponential convergence under a spectral gap}\label{sec:ntkgf:convergence}

Throughout this section assume $m\ge1$, and say that $K$ satisfies the \emph{Rayleigh--Ritz condition with parameter $\lambda$} if
$\lambda\|v\|^2\le v^\top Kv$ for all $v\in\mathbb R^m$.

\begin{theorem}[Exponential decay of the residual and the loss]\label{thm:ntkgf:expDecay}
  \uses{lem:ntkf:gronwallDecay, thm:ntkgf:outputODE}
  \lean{NTK.dotProduct_mulVec_nonneg_of_posSemidef, NTK.inner_toLp_mulVec_eq_dotProduct,
        NTK.deriv_norm_sq_timeVarying_ode, NTK.deriv_norm_sq_linear_ode,
        NTK.deriv_norm_sq_le_of_rayleighRitz_timeVarying, NTK.deriv_norm_sq_le_of_rayleighRitz,
        NTK.residual_norm_sq_exponential_decay_timeVarying,
        NTK.residual_norm_exponential_decay_timeVarying,
        NTK.mse_loss_exponential_decay_timeVarying,
        NTK.residual_norm_sq_exponential_decay, NTK.residual_norm_exponential_decay,
        NTK.mse_loss_exponential_decay}
  \leanok
Let $r:\mathbb R\to\mathbb R^m$ satisfy $\partial_tr(t)=-\frac1mK(t)r(t)$ for all $t$, where each $K(t)$ satisfies the Rayleigh--Ritz condition
with parameter $\lambda$ (time-varying kernel), or $K(t)\equiv K_\infty$ (autonomous). Then for all $t\ge0$
\[
  \|r(t)\|\le\|r(0)\|\,e^{-\lambda t/m},\qquad
  \tfrac1{2m}\|r(t)\|^2\le\tfrac1{2m}\|r(0)\|^2\,e^{-2\lambda t/m}.
\]
\end{theorem}

\begin{proof}
  \leanok
\emph{Step 1.} $\frac{d}{dt}\|r(t)\|^2=2r^\top\partial_tr=-\frac2m\,r(t)^\top K(t)r(t)$.

\emph{Step 2.} By the Rayleigh--Ritz condition this is at most $-\frac{2\lambda}{m}\|r(t)\|^2$.

\emph{Step 3.} Gr\"onwall's inequality (Lemma~\ref{lem:ntkf:gronwallDecay}) gives $\|r(t)\|^2\le\|r(0)\|^2e^{-2\lambda t/m}$, and taking square roots gives the first bound.

\emph{Step 4.} Multiplying by $\frac1{2m}$ gives the loss bound.
\end{proof}

\begin{theorem}[Local-in-time version]\label{thm:ntkgf:expDecayIcc}
  \uses{thm:ntkgf:expDecay, lem:ntkf:gronwallDecay}
  \lean{NTK.residual_norm_sq_exponential_decay_timeVarying_Icc,
        NTK.residual_norm_exponential_decay_timeVarying_Icc}
  \leanok
Let $T\ge0$, let $r$ be continuous on $[0,T]$ with $\partial_tr(t)=-\frac1mK(t)r(t)$ for $t\in(0,T)$ only, and suppose $K(s)$ satisfies the
Rayleigh--Ritz condition with parameter $\lambda$ for all $s\in[0,T]$. Then for $t\in[0,T]$,
$\|r(t)\|^2\le\|r(0)\|^2e^{-2\lambda t/m}$ and $\|r(t)\|\le\|r(0)\|e^{-\lambda t/m}$.
\end{theorem}

\begin{proof}
  \uses{lem:ntkf:gronwallDecay}
  \leanok
Same computation as above, using the interval version of Gr\"onwall's inequality that needs differentiability only on $(0,T)$; this is what allows
the estimate to be used inside a bootstrap, where the spectral gap is known only while the trajectory remains in a ball.
\end{proof}

\begin{lemma}[Loss decay implies convergence]\label{lem:ntkgf:tendstoZero}
  \lean{NTK.tendsto_zero_of_le_mul_exp_neg}
  \leanok
If $0\le L(t)\le L_0e^{-ct}$ for $t\ge0$ and $c>0$, then $L(t)\to0$ as $t\to\infty$.
\end{lemma}

\begin{proof}
  \leanok
$e^{-ct}\to0$ and squeeze.
\end{proof}


\section{Closed-form linear dynamics}\label{sec:ntkgf:linear}

In the infinite-width limit the kernel is constant, $K_t\equiv K_\infty$, and the residual obeys the linear ODE $\partial_tr=-\frac1mK_\infty r$.

\begin{theorem}[Matrix-exponential solution]\label{thm:ntkgf:matrixExp}
  \uses{thm:ntkgf:expDecay}
  \lean{NTK.matrix_exp_residual_trajectory_zero, NTK.matrix_exp_output_trajectory_zero,
        NTK.matrix_exp_residual_eq_output_sub_y, NTK.hasDerivAt_clm_apply_const,
        NTK.toEuclideanVecCLM, NTK.matrix_exp_residual_trajectory_hasDerivAt,
        NTK.matrix_exp_output_trajectory_hasDerivAt, NTK.matrix_exp_residual_decay,
        NTK.matrix_exp_loss_decay}
  \leanok
For a matrix $K_\infty\in\mathbb R^{m\times m}$ and $r_0,f_0,y\in\mathbb R^m$ define
\[
  r(t)=\exp\!\bigl(-\tfrac tmK_\infty\bigr)r_0,\qquad f(t)=y+\exp\!\bigl(-\tfrac tmK_\infty\bigr)(f_0-y).
\]
Then $r(0)=r_0$, $f(0)=f_0$, $f(t)-y=r(t)$, and $\partial_tr=-\frac1mK_\infty r$, $\partial_tf=-\frac1mK_\infty(f-y)$. 

If $K_\infty$ satisfies the
Rayleigh--Ritz condition with parameter $\lambda$, then for $t\ge0$,
$\|r(t)\|\le\|r_0\|e^{-\lambda t/m}$ and $\frac1{2m}\|r(t)\|^2\le\frac1{2m}\|r_0\|^2e^{-2\lambda t/m}$.
\end{theorem}

\begin{proof}
  \leanok
The matrix exponential satisfies $\frac{d}{ds}\exp(sA)=A\exp(sA)$, so with $A=-K_\infty/m$ the curve $t\mapsto\exp(tA)r_0$ solves the ODE; evaluation at a fixed
vector is a continuous linear map and therefore commutes with differentiation. The decay estimates follow from Theorem~\ref{thm:ntkgf:expDecay}
applied to this solution.
\end{proof}

\begin{theorem}[Eigenmode decomposition of the frozen residual]\label{thm:ntkgf:eigenmodes}
  \uses{thm:ntkgf:matrixExp}
  \lean{NTK.matrix_exp_smul_mulVec_of_eigenvector, NTK.inner_matrix_exp_mulVec_of_eigenvector,
        NTK.inner_eigenvectorBasis_matrix_exp_mulVec, NTK.matrix_exp_mulVec_eq_sum_eigenmodes,
        NTK.norm_sq_matrix_exp_mulVec_eq_sum}
  \leanok
If $Kv=\lambda v$ then $\exp(sK)v=e^{s\lambda}v$ (no symmetry needed). 

If $K$ is symmetric, then $\langle v,\exp(sK)r\rangle=e^{s\lambda}\langle v,r\rangle$.

For a Hermitian $K$ with orthonormal eigenbasis $(v_k)$ and eigenvalues $(\lambda_k)$ (Mathlib's \texttt{eigenvectorBasis}),
\[
  \langle v_k,\exp(sK)r\rangle=e^{s\lambda_k}\langle v_k,r\rangle,\quad
  \exp(sK)r=\sum_ke^{s\lambda_k}\langle v_k,r\rangle\,v_k,\quad
  \|\exp(sK)r\|^2=\sum_ke^{2s\lambda_k}\langle v_k,r\rangle^2 .
\]
With $s=-t/m$ this is the decoupled mode equation $r_k(t)=e^{-\lambda_kt/m}r_k(0)$ of the frozen-kernel residual, and Parseval for its energy.
\end{theorem}

\begin{proof}
  \leanok
$\exp(sK)v=\sum_n\frac{s^n}{n!}K^nv=\sum_n\frac{(s\lambda)^n}{n!}v=e^{s\lambda}v$. For symmetric $K$, $\exp(sK)^\top=\exp(sK)$, so
$\langle v,\exp(sK)r\rangle=\langle\exp(sK)v,r\rangle$. The expansion in the eigenbasis is the orthonormal-basis expansion
$w=\sum_k\langle v_k,w\rangle v_k$ with $w=\exp(sK)r$, and Parseval gives the energy.
\end{proof}

\begin{remark}[Why the eigenmode decomposition is useful]\label{rmk:ntkgf:eigenmodesUse}
  \uses{thm:ntkgf:eigenmodes, thm:ntkgf:matrixExp}
The frozen-kernel residual $r(t)=\exp(-\tfrac tmK_\infty)r_0$ solves a system of $m$ \emph{coupled} equations, one for each training point. Writing it in the eigenbasis of $K_\infty$ decouples the system into $m$ independent scalar equations:
the component of the residual along the eigenvector $v_k$ is $r_k(t)=e^{-\lambda_kt/m}\,r_k(0)$ and simply decays at its own rate $\lambda_k/m$. This makes the training dynamics easy to read off.
\begin{itemize}
  \item \emph{Which directions are learned fast or slowly.} Components of the initial residual along eigenvectors with a large eigenvalue $\lambda_k$ are fitted quickly; those with a small eigenvalue are fitted slowly, and components along an eigenvalue $0$ are never fitted. The slowest rate $\lambda_{\min}/m$ is the bottleneck, which is exactly the rate in the exponential decay bound of Theorem~\ref{thm:ntkgf:expDecay}; the eigenmode formula shows that this bound is attained when the initial residual lies in the slowest mode.
  \item \emph{The loss is an explicit sum.} By Parseval, $L(t)=\tfrac1{2m}\|r(t)\|^2=\tfrac1{2m}\sum_ke^{-2\lambda_kt/m}\langle v_k,r_0\rangle^2$, so one sees exactly how much of the initial error remains at time $t$ and in which modes.
  \item \emph{It is the reference curve for the real, non-frozen dynamics.} The actual residual is compared with these modes in Theorem~\ref{thm:ntkgf:residualFrozen}(3): each mode of the true residual stays within $\tfrac1m\varepsilon_K\|r_0\|\,t$ of its frozen curve, which is how the lazy-training results of the later chapters turn ``the kernel barely moves'' into ``the network follows the frozen-kernel dynamics''.
\end{itemize}
\end{remark}

\begin{remark}[Role of this section in the NTK development]\label{rmk:ntkgf:linearRole}
  \uses{thm:ntkgf:matrixExp, thm:ntkgf:eigenmodes, thm:ntkgf:outputODE, thm:ntkgf:expDecay}
\emph{Main result.} If the kernel is frozen at a fixed matrix $K_\infty$, the training dynamics can be solved exactly. The residual is $r(t)=\exp(-\tfrac tmK_\infty)r_0$ and the outputs are $f(t)=y+\exp(-\tfrac tmK_\infty)(f_0-y)$ (Theorem~\ref{thm:ntkgf:matrixExp})

In the eigenbasis of $K_\infty$ each mode decays independently at the rate $\lambda_k/m$ (Theorem~\ref{thm:ntkgf:eigenmodes}). This is the \emph{linear} (frozen-kernel) model of training, and the NTK picture says that a very wide network behaves like it.

\emph{What it builds on.} The equation solved here, $\partial_tr=-\frac1mK_\infty r$, is the residual ODE of Theorem~\ref{thm:ntkgf:outputODE} with the time-dependent kernel $K(\theta(t))$ replaced by a constant. The exponential decay estimate of Theorem~\ref{thm:ntkgf:expDecay}, which holds for any time-varying kernel with a spectral gap, is only an \emph{inequality}; for a constant kernel the matrix exponential gives the exact trajectory, and the decay bound is recovered as a corollary (the last statement of Theorem~\ref{thm:ntkgf:matrixExp}).

\emph{What it is used for.} The section itself makes no claim about a real network, since for a finite network the kernel is never exactly constant. Its role is to be the reference curve against which the true dynamics are compared.
\begin{itemize}
  \item The next section (Theorem~\ref{thm:ntkgf:residualFrozen}) bounds the distance between the true residual and $\exp(-\frac tmK_\infty)r_0$, and between their eigenmode components, by $\frac1m\varepsilon_K\|r_0\|t$, where $\varepsilon_K$ bounds how far the kernel moves.
  \item Kernel stability and lazy training (Sections~\ref{sec:ntkgf:kernelStability} and~\ref{sec:ntkgf:bootstrap}) are what supply a small $\varepsilon_K$, by keeping the parameters near $\theta_0$.
  \item The affine flow (Theorem~\ref{thm:ntkgf:affineFlow}) is solved with the same matrix exponential, now with $K_\infty=JJ^\top$ for a fixed Jacobian $J$, and its limit is kernel regression (Theorem~\ref{thm:ntkgf:affineLimit}).
  \item For test points, the frozen-kernel prediction $-a^\top\bigl(r_0-\exp(-\frac tmK_\infty)r_0\bigr)$ is compared with the actual prediction in Theorems~\ref{thm:ntkgf:predictionError} and~\ref{thm:ntkgf:predictionUniform}.
  \item In the two-layer chapter, the random initial residual $r_n(0)$ converges in distribution to a Gaussian $G-y$ and $K_n(0)\to K_\infty$, so the trained residual converges to $\exp(-\frac tmK_\infty)(G-y)$ (Theorem~\ref{thm:ntktl:residualLimit}).
\end{itemize}
\end{remark}


\section{Stability of linear ODEs under coefficient perturbation}\label{sec:ntkgf:odeStability}

\begin{theorem}[Perturbation of a dissipative linear ODE]\label{thm:ntkgf:odePerturbation}
  \lean{NTK.norm_sub_le_of_linear_ode_perturbation}
  \leanok
Let $A,B:\mathbb R\to\mathbb R^{m\times m}$, let $r$ be continuous on $[0,T]$ with $r'=-A(t)r$ on $(0,T)$, let $s'=-B(t)s$ on $[0,T]$, and assume
$A(t)\succeq0$ for $t\in[0,T]$ and $\|A(t)-B(t)\|\,\|s(t)\|\le a$ on $[0,T]$. Then
\[
  \|r(t)-s(t)\|\le\|r(0)-s(0)\|+a\,t,\qquad t\in[0,T].
\]
\end{theorem}

\begin{proof}
  \leanok
Let $e=r-s$. Then $e'=-A(t)e-(A(t)-B(t))s$. Hence $\frac{d}{dt}\tfrac12\|e\|^2=-\langle e,A(t)e\rangle-\langle e,(A-B)s\rangle\le\|e\|\,a$, because $A(t)\succeq0$ makes the
dissipative term nonpositive. Dividing by $\|e\|$ (or arguing with $\sqrt{\|e\|^2+\varepsilon}$) gives $\frac d{dt}\|e\|\le a$ and integration gives the claim. The
positive semidefiniteness of $A$ is what removes any exponential Gr\"onwall factor.
\end{proof}

\begin{theorem}[Residual versus frozen-kernel residual]\label{thm:ntkgf:residualFrozen}
  \uses{thm:ntkgf:odePerturbation, thm:ntkgf:matrixExp, thm:ntkgf:eigenmodes}
  \lean{NTK.residual_sub_frozen_residual_le, NTK.residual_sub_matrix_exp_le,
        NTK.abs_inner_eigenvector_residual_sub_mode_le}
  \leanok
Let $r$ be continuous on $[0,T]$ with $\partial_tr=-\frac1mK(t)r$ on $(0,T)$, where $K(t)\succeq0$ and $K_\infty\succeq0$.
\begin{enumerate}
  \item If $s'=-\frac1mK_\infty s$ and $\|K(t)-K_\infty\|\,\|s(t)\|\le b$ on $[0,T]$, then $\|r(t)-s(t)\|\le\|r(0)-s(0)\|+\frac1mbt$.
  \item If $\|K(t)-K_\infty\|\le\varepsilon_K$ on $[0,T]$, then for $t\in[0,T]$
        $\bigl\|r(t)-\exp(-\tfrac tmK_\infty)r(0)\bigr\|\le\frac1m\,\varepsilon_K\,\|r(0)\|\,t.$
  \item Under the hypotheses of (2), for every eigenpair $(\lambda_k,v_k)$ of $K_\infty$ (with orthonormal $v_k$),
        $\bigl|\langle v_k,r(t)\rangle-e^{-\lambda_kt/m}\langle v_k,r(0)\rangle\bigr|\le\frac1m\varepsilon_K\|r(0)\|\,t$.
\end{enumerate}
\end{theorem}

\begin{proof}
  \uses{thm:ntkgf:odePerturbation, thm:ntkgf:matrixExp, thm:ntkgf:eigenmodes}
  \leanok
(1) is Theorem~\ref{thm:ntkgf:odePerturbation} with $A=K/m$, $B=K_\infty/m$ and $a=b/m$. For (2) take $s(t)=\exp(-\frac tmK_\infty)r(0)$, which starts at $r(0)$, and note
$\|s(t)\|\le\|r(0)\|$ because $K_\infty\succeq0$; thus $\|K-K_\infty\|\|s\|\le\varepsilon_K\|r(0)\|$. For (3), $|\langle v_k,r-s\rangle|\le\|r-s\|$ and
$\langle v_k,s(t)\rangle=e^{-\lambda_kt/m}\langle v_k,r(0)\rangle$ by Theorem~\ref{thm:ntkgf:eigenmodes}.
\end{proof}

\begin{remark}[Role of this section in the NTK development]\label{rmk:ntkgf:odeStabilityRole}
  \uses{thm:ntkgf:odePerturbation, thm:ntkgf:residualFrozen, thm:ntkgf:matrixExp, thm:ntkgf:eigenmodes, thm:ntkgf:outputODE}
\emph{Main result.} The true residual obeys $\partial_tr=-\frac1mK(\theta(t))\,r$ with a kernel that moves, whereas the previous section solved the same equation exactly for a constant kernel $K_\infty$. This section shows that the two solutions stay close when the kernel is nearly constant. Theorem~\ref{thm:ntkgf:odePerturbation} is the abstract statement: two linear ODEs whose coefficient matrices differ by a bounded amount produce solutions that differ by at most the initial gap plus $a\,t$, \emph{provided the true coefficient is positive semidefinite}. Theorem~\ref{thm:ntkgf:residualFrozen} specializes it to $A=K/m$, $B=K_\infty/m$. If $\|K(t)-K_\infty\|\le\varepsilon_K$ on $[0,T]$, then
\[
  \bigl\|r(t)-\exp(-\tfrac tmK_\infty)r(0)\bigr\|\le\tfrac1m\,\varepsilon_K\,\|r(0)\|\,t,
\]
and every eigenmode of the true residual stays within the same distance of its frozen curve.

\emph{What it builds on.}
\begin{itemize}
  \item The equation being perturbed is the residual ODE of Theorem~\ref{thm:ntkgf:outputODE}, and the positive semidefiniteness of $K$ that makes the estimate free of an exponential factor is Lemma~\ref{lem:ntkgf:ntkPsd}.
  \item The reference solution $\exp(-\frac tmK_\infty)r(0)$ and its eigenmode components are those of Theorems~\ref{thm:ntkgf:matrixExp} and~\ref{thm:ntkgf:eigenmodes}. The bound $\|s(t)\|\le\|r(0)\|$ used in part~(2) of Theorem~\ref{thm:ntkgf:residualFrozen} is the decay already established there.
\end{itemize}

\emph{Why no Gr\"onwall factor matters.} A general perturbation estimate for linear ODEs would grow like $e^{\|A\|t}$, which is useless on the time scales of interest because $\|K\|/m$ is not small. Dissipativity replaces that factor by the linear bound $a\,t$, so the error is small as soon as $\varepsilon_K\,t/m$ is small. This is why the kernel only has to be nearly constant on a time interval of fixed length, and no smallness relative to $\|K\|$ is needed.

\emph{What it is used for.} The theorem takes $\varepsilon_K$ as a hypothesis and says nothing about why the kernel should be nearly constant. That input comes later.
\begin{itemize}
  \item Kernel stability and lazy training (Sections~\ref{sec:ntkgf:kernelStability} and~\ref{sec:ntkgf:bootstrap}) show that $\|K(\theta)-K(\theta_0)\|$ is small whenever $\theta$ stays near $\theta_0$, which supplies $\varepsilon_K$ with $K_\infty=K(\theta_0)$ or its limit.
  \item The test-point prediction error (Theorems~\ref{thm:ntkgf:predictionError} and~\ref{thm:ntkgf:predictionUniform}) integrates the bound of part~(2) to compare the true prediction at a new input with the frozen-kernel prediction.
  \item In the two-layer chapter, Theorem~\ref{thm:ntktl:residualLimit} combines the bound with kernel stationarity as the width tends to infinity, to show that the trained residual is asymptotically $\exp(-\frac tmK_\infty)(G-y)$.
\end{itemize}
\end{remark}


\section{Kernel stability: Lipschitz propagation and Rayleigh quotients}\label{sec:ntkgf:kernelStability}

Recall (Definition~\ref{def:ntkgf:jacobian}) that $J(\theta)\in\mathbb R^{m\times P}$ is the output Jacobian, whose $\alpha$-th row is the gradient $\nabla_\theta f(x^\alpha;\theta)^\top$ of the output at training point $\alpha$ with respect to the $P$ parameters, and that $K(\theta)=J(\theta)J(\theta)^\top\in\mathbb R^{m\times m}$ is the empirical NTK matrix, the Gram matrix of those gradients.
So $K$ is a quadratic function of $J$, and the theorem says that a bound on $J$ and on how fast $J$ changes with $\theta$ controls how fast the kernel $K$ changes. Norms of matrices are Frobenius norms.

\begin{theorem}[Lipschitz propagation from the Jacobian to the kernel]\label{thm:ntkgf:kernelLipschitz}
  \uses{def:ntkgf:jacobian, lem:ntkf:mulVec}
  \lean{NTK.matrix_mul_transpose_sub_mul_transpose, NTK.norm_mul_transpose_sub_mul_transpose_le,
        NTK.empiricalNTKMatrix_sub_empiricalNTKMatrix, NTK.empiricalNTKMatrix_sub_le_of_jacobian_bound,
        NTK.empiricalNTKMatrix_sub_le_of_jacobian_lipschitz,
        NTK.empiricalNTKMatrix_lipschitz_of_jacobian_bound}
  \leanok

For matrices $A,B\in\mathbb R^{m\times P}$, $AA^\top-BB^\top=(A-B)A^\top+B(A-B)^\top$, hence if $\|A\|,\|B\|\le M$ then
$\|AA^\top-BB^\top\|_{\mathrm F}\le2M\|A-B\|_{\mathrm F}$. Consequently, if $\|J(\theta_1)\|,\|J(\theta_2)\|\le M$,
\[
  \|K(\theta_1)-K(\theta_2)\|\le2M\|J(\theta_1)-J(\theta_2)\|,
\]
and if moreover $\|J(\theta_1)-J(\theta_2)\|\le L_J\|\theta_1-\theta_2\|$, then $\|K(\theta_1)-K(\theta_2)\|\le2ML_J\|\theta_1-\theta_2\|$.
In particular, on any set $S\ni\theta_0$ on which $\|J(\theta)\|\le M$ and $\|J(\theta)-J(\theta_0)\|\le L_J\|\theta-\theta_0\|$,
\[
  \|K(\theta)-K(\theta_0)\|\le(2ML_J)\,\|\theta-\theta_0\|\qquad(\theta\in S).
\]
\end{theorem}

\begin{proof}
  \leanok
The algebraic decomposition is a direct expansion. Apply the triangle inequality and $\|XY\|_{\mathrm F}\le\|X\|_{\mathrm F}\|Y\|_{\mathrm F}$, using $\|A^\top\|=\|A\|$.
\end{proof}

\begin{theorem}[$C^{1,1}$ Taylor bound]\label{thm:ntkgf:taylor}
  \lean{NTK.norm_sub_sub_fderiv_le_of_lipschitz_fderiv, NTK.norm_trainingOutputs_sub_linearization_le}
  \leanok
Let $G:E\to F$ be Fr\'echet differentiable on the closed ball of radius $r$ about $x_0$ with derivative $G'(z)$, and suppose
$\|G'(z)-G'(x_0)\|\le L\|z-x_0\|$ there. Then for $\|x-x_0\|\le r$,
\[
  \|G(x)-G(x_0)-G'(x_0)(x-x_0)\|\le\frac L2\|x-x_0\|^2 .
\]
For the training outputs: if the output Jacobian satisfies $\|J(\theta)-J(\theta_0)\|\le L\|\theta-\theta_0\|$ on that ball and each output is differentiable, then
$\|f(\theta)-f(\theta_0)-J(\theta_0)(\theta-\theta_0)\|\le\frac L2\|\theta-\theta_0\|^2$.
\end{theorem}

\begin{proof}
  \leanok
Let $\phi(s)=G(x_0+s(x-x_0))-G(x_0)-sG'(x_0)(x-x_0)$ for $s\in[0,1]$. Then $\phi'(s)=(G'(x_0+s(x-x_0))-G'(x_0))(x-x_0)$ has norm at most $Ls\|x-x_0\|^2$. The mean-value
inequality along $[0,1]$ gives $\|\phi(1)\|\le\int_0^1Ls\|x-x_0\|^2ds=\frac L2\|x-x_0\|^2$.

\emph{Comparison with Mathlib's Taylor theorems.} The first-order remainder estimate in Mathlib (\texttt{taylor\_mean\_remainder\_bound} with $n=1$) assumes that the function is $C^2$ on the interval, i.e.\ twice continuously differentiable, together with a bound $C'$ on the norm of its second derivative; it then gives $\|f(x)-f(a)-f'(a)(x-a)\|\le C'(x-a)^2$, which is off by a factor $2$ from the sharp constant $\tfrac12$.
The Lagrange form (\texttt{taylor\_mean\_remainder\_lagrange}) has the sharp constant, but only for real-valued functions, and it expresses the remainder through the second derivative at an unspecified intermediate point, so it also needs that derivative to exist. Both are statements about functions of one real variable, so applying them here would also require restricting to the segment from $x_0$ to $x$, as above.
Our hypothesis is weaker: it only asks that $G'$ is $L$-Lipschitz, and $G$ need not be twice differentiable (for example $G'(s)=|s|$ is Lipschitz but not differentiable at $0$). If $G''$ exists with $\|G''\|\le L$ then $G'$ is $L$-Lipschitz, so this is the more general condition; it also covers maps between normed spaces and keeps the sharp constant $\tfrac L2$.
For the training outputs take $G=f(\cdot)$ and $G'(z)v=J(z)v$, whose operator norm is at most $\|J(z)\|_{\mathrm F}$.
\end{proof}

\begin{theorem}[Rayleigh quotient stability]\label{thm:ntkgf:rayleighStability}
  \uses{thm:ntkgf:kernelLipschitz}
  \lean{NTK.abs_dotProduct_mulVec_sub_le, NTK.rayleigh_lower_bound_of_sub_smul_posSemidef,
        NTK.rayleigh_quotient_lower_bound_of_matrix_dist,
        NTK.rayleigh_quotient_lower_bound_of_displacement}
  \leanok
\begin{enumerate}
  \item For matrices $A,B$ and vectors $v$, $|v^\top Av-v^\top Bv|\le\|A-B\|\,\|v\|^2$.
  \item If $K-\lambda I\succeq0$ then $\lambda\|v\|^2\le v^\top Kv$ for all $v$.
  \item If $\lambda_0\|v\|^2\le v^\top K_0v$ and $\|K-K_0\|\le\varepsilon$, then $(\lambda_0-\varepsilon)\|v\|^2\le v^\top Kv$.
  \item If the empirical NTK matrix at $\theta_0$ satisfies the Rayleigh--Ritz condition with parameter $\lambda_0$, the output Jacobian is $M$-bounded at $\theta$ and $\theta_0$ and
        $\|J(\theta)-J(\theta_0)\|\le L_J\|\theta-\theta_0\|$, then the empirical NTK matrix at $\theta$ satisfies it with parameter
        $\lambda_0-2ML_J\|\theta-\theta_0\|$.
\end{enumerate}
\end{theorem}

\begin{proof}
  \leanok
(1) $|v^\top(A-B)v|\le\|(A-B)v\|\|v\|\le\|A-B\|\|v\|^2$ (Lemma~\ref{lem:ntkf:mulVec}). (2) $v^\top Kv=v^\top(K-\lambda I)v+\lambda\|v\|^2\ge\lambda\|v\|^2$. (3) follows from (1).
(4) is (3) with $\varepsilon=2ML_J\|\theta-\theta_0\|$ from Theorem~\ref{thm:ntkgf:kernelLipschitz}.
\end{proof}

\begin{remark}[Why Rayleigh quotient stability matters]\label{rmk:ntkgf:rayleighSignificance}
  \uses{thm:ntkgf:rayleighStability, thm:ntkgf:expDecayIcc, thm:ntkgf:rayleighBall}
The exponential convergence of the loss (Theorems~\ref{thm:ntkgf:expDecay} and~\ref{thm:ntkgf:expDecayIcc}) needs a spectral gap of the empirical NTK matrix $K(\theta(t))$ at \emph{every} parameter $\theta(t)$ that the flow visits. A gap is only known at the initialization $\theta_0$, where the initial kernel concentrates around its limit.

Part~(4) of the theorem bridges this: if the kernel has gap $\lambda_0$ at $\theta_0$, then at any $\theta$ with $2ML_J\|\theta-\theta_0\|<\lambda_0$ it still has the positive gap $\lambda_0-2ML_J\|\theta-\theta_0\|$. In other words, a spectral gap is robust: it survives, with a smaller constant, as long as the parameters stay close enough to $\theta_0$.

This is the ingredient that makes the lazy-training bootstrap work (Theorems~\ref{thm:ntkgf:rayleighBall} and~\ref{thm:ntkgf:lazyDisplacement}): the parameters stay near $\theta_0$ because the residual decays, and the residual decays because the kernel keeps a gap near $\theta_0$; the stability estimate is what lets this circular argument be closed by continuous induction.
\end{remark}


\section{Lazy training: displacement bound and the bootstrap}\label{sec:ntkgf:bootstrap}

The central point is a circularity. By \emph{Jacobian estimates} we mean two quantitative assumptions on the output Jacobian $J(\theta)$ (Definition~\ref{def:ntkgf:jacobian}), the matrix of derivatives of the training outputs with respect to the parameters:
\begin{itemize}
  \item \emph{boundedness}: $\|J(\theta)\|\le M$, so the outputs cannot change faster than a fixed rate $M$ when the parameters move;
  \item \emph{Lipschitz continuity at $\theta_0$}: $\|J(\theta)-J(\theta_0)\|\le L_J\|\theta-\theta_0\|$, so the Jacobian, and with it the kernel $K=JJ^\top$ (Theorem~\ref{thm:ntkgf:kernelLipschitz}), stays close to its initial value.
\end{itemize}
Together they give the spectral gap of $K(\theta)$ needed for exponential decay of the residual (Theorem~\ref{thm:ntkgf:rayleighStability}). For a wide network one can verify them at initialization only on a ball $\{\|\theta-\theta_0\|\le r\}$, with $r$ small enough that $2ML_Jr\le\lambda_0/2$. But the conclusion we want is that the parameters never leave this ball (small displacement), and to use the estimates at time $t$ we must already know $\theta(t)$ is inside the ball.
The bootstrap principle (Theorem~\ref{thm:ntkf:bootstrap}) resolves this by continuous induction on the time $S$ up to which the trajectory is known to stay in the ball. Below, $\theta$ is a forward gradient flow of the MSE loss
(Definition~\ref{def:ntkgf:forwardGF}) and every output is differentiable along it.

\begin{theorem}[Displacement of a linear image of a gradient flow]\label{thm:ntkgf:displacementLinear}
  \uses{def:ntkgf:forwardGF, def:ntkgf:outputs, def:ntkgf:jacobian, thm:ntkgf:mseGradient}
  \lean{NTK.restrictCoords, NTK.restrictCoords_apply, NTK.norm_sq_restrictCoords,
        NTK.sum_sq_mulVec_transpose_le, NTK.norm_restrictCoords_gradient_mseLoss_le,
        NTK.norm_map_sub_le_integral_of_forwardGF}
  \leanok
Let $\theta$ be a forward gradient flow (Definition~\ref{def:ntkgf:forwardGF}) of a loss $L$ on $\mathbb R^P$, let $F$ be a normed space and $\mathrm{Lin}:\mathbb R^P\to F$ a continuous linear map.
\begin{enumerate}
  \item If $\|\mathrm{Lin}(\nabla L(\theta(t)))\|_F\le B(t)$ on $[0,T]$ with $B$ interval integrable, then
        $\|\mathrm{Lin}(\theta(T)-\theta_0)\|_F\le\int_0^TB(t)\,dt$.
  \item Let $\kappa$ be a finite index set, $k:\kappa\to\{1,\dots,P\}$, and $\mathrm{restrictCoords}\,k:\mathbb R^P\to\mathbb R^\kappa$, $v\mapsto(v_{k(o)})_{o\in\kappa}$. For the MSE loss $L$ with residual $r$ (Definition~\ref{def:ntkgf:outputs}) and output Jacobian $J$ (Definition~\ref{def:ntkgf:jacobian}),
        \[
          \bigl\|\mathrm{restrictCoords}\,k\,(\nabla L(\theta))\bigr\|\le\frac1m\Bigl(\sum_{\alpha}\sum_{o\in\kappa}J_{\alpha,k(o)}(\theta)^2\Bigr)^{1/2}\|r(\theta)\| .
        \]
\end{enumerate}
\end{theorem}

\begin{proof}
  \uses{def:ntkgf:forwardGF, thm:ntkgf:mseGradient}
  \leanok
Since $\theta'=-\nabla L(\theta)$ (Definition~\ref{def:ntkgf:forwardGF}) and $\mathrm{Lin}$ is linear and continuous, $\mathrm{Lin}(\theta(T)-\theta_0)=-\int_0^T\mathrm{Lin}(\nabla L(\theta(t)))dt$, whose norm is at most $\int_0^TB$. Only the ODE for $t>0$ and continuity on
$[0,\infty)$ are used. For the block bound, $\nabla L=\frac1mJ^\top r$ (Theorem~\ref{thm:ntkgf:mseGradient}) and Cauchy--Schwarz gives $\sum_o(\sum_\alpha r^\alpha J_{\alpha o})^2\le\|r\|^2\sum_{\alpha,o}J_{\alpha o}^2$.
\end{proof}

\begin{remark}[Reading the statement: $\mathrm{Lin}$, blocks of coordinates, and the speed of a block]\label{rmk:ntkgf:displacementLinearReading}
  \uses{thm:ntkgf:displacementLinear}
The loss may more generally live on a real Hilbert space $E$ in place of $\mathbb R^P$. The map $\mathrm{Lin}$ extracts or summarizes some feature of the parameter vector. Two cases are used: $F=\mathbb R^P$ with $\mathrm{Lin}=\mathrm{id}$, which gives the displacement of the whole parameter vector, and $F=\mathbb R^\kappa$ with $\mathrm{Lin}=\mathrm{restrictCoords}\,k$, which gives the displacement of a subset of the parameters, for example the weights of a single neuron.

A \emph{block of coordinates} is a chosen subset of the $P$ parameters, regarded as one group: the sub-vector of $\theta$ made of those entries only. For the two-layer network, a natural block is all the parameters belonging to one hidden neuron $i$, namely its input weights $W_{i1},\dots,W_{id}$ and its readout weight $a_i$, so $d+1$ of the $nd+n$ coordinates (the parametrization by a flat parameter vector is Definition~\ref{def:ntktl:packing}).
In part~(2), $\kappa$ is the finite set of labels of the chosen coordinates, and $k$ sends each label $o\in\kappa$ to the number of the parameter coordinate it selects. In Lean $k$ is an arbitrary function; in the intended use it is injective, so each coordinate is selected once.

The \emph{speed of the block} is how fast those parameters move, namely the norm of the corresponding entries of $\theta'=-\nabla L(\theta)$; part~(2) bounds it. Part~(1) then bounds how far the block travels in time $T$ by integrating its speed. Since this speed is smaller than that of the whole parameter vector, the displacement of each neuron can be controlled by its own size, as used in Theorem~\ref{thm:ntktl:neuronDisplacement}.
\end{remark}

\begin{lemma}[Loss monotonicity and continuity of the residual]\label{lem:ntkgf:lossMonotone}
  \lean{NTK.mseLoss_le_of_hasDerivWithinAt_neg_gradient, NTK.continuousOn_trainingResidual_comp}
  \leanok
If $\theta$ solves $\theta'=-\nabla L(\theta)$ on $[0,S]$ (one-sided at the endpoints) with differentiable outputs, then $L(\theta(t))\le L(\theta(0))$ for $t\in[0,S]$.
The residual along a curve is continuous on any set where the curve is continuous and every output is differentiable at the curve.
\end{lemma}

\begin{proof}
  \leanok
$\frac d{dt}L(\theta(t))=-\|\nabla L(\theta(t))\|^2\le0$; this needs neither a global trajectory nor differentiability of $L$ off the curve. Continuity of the residual is the continuity of a differentiable composition.
\end{proof}

\begin{remark}[The residual along a curve]\label{rmk:ntkgf:residualAlongCurve}
  \uses{lem:ntkgf:lossMonotone, def:ntkgf:outputs}
Recall from Definition~\ref{def:ntkgf:outputs} that the residual at a parameter $\theta\in\mathbb R^P$ is the vector of prediction errors on the training set,
$r(\theta)=f(\theta)-y\in\mathbb R^m$, with entries $r^\alpha(\theta)=f(x^\alpha;\theta)-y^\alpha$.
Given a curve $t\mapsto\theta(t)$ in parameter space, such as a gradient flow, the \emph{residual along the curve} is the function of time $t\mapsto r(\theta(t))\in\mathbb R^m$, i.e.\ the training errors of the network as it is being trained.
The second statement of the lemma says this function is continuous at those times $t$ where the curve $\theta$ is continuous and each of the $m$ outputs $\theta\mapsto f(x^\alpha;\theta)$ is differentiable at $\theta(t)$ (differentiability implies continuity there, so each entry $r^\alpha(\theta(t))$ is a composition of continuous maps).
Continuity of $t\mapsto r(\theta(t))$, hence of $\|r(\theta(t))\|$, is what allows the residual and the loss to be handled as ordinary continuous functions of time in the bootstrap arguments below and in the two-layer chapter.
\end{remark}

\begin{theorem}[Displacement bounds along the flow]\label{thm:ntkgf:displacementBound}
  \uses{thm:ntkgf:displacementLinear, thm:ntkgf:expDecayIcc, thm:ntkgf:mseGradient}
  \lean{NTK.displacement_le_integral_of_rayleigh, NTK.displacement_integral_bound,
        NTK.displacement_bound_of_psd}
  \leanok
Let $\theta$ be a forward gradient flow of the MSE loss with $\|J(\theta(t))\|\le M$ for $t\in[0,T]$.
\begin{enumerate}
  \item If the empirical NTK satisfies the Rayleigh--Ritz condition with parameter $\lambda$ (any real, e.g.\ $0$) at every $\theta(t)$, $t\in[0,T]$, then
        $\|\theta(T)-\theta_0\|\le\int_0^T\frac Mm\|r_0\|e^{-\lambda t/m}dt$, where $r_0=r(\theta_0)$.
  \item If $\lambda>0$, then $\|\theta(T)-\theta_0\|\le M\|r_0\|/\lambda$, a bound with no dependence on $T$.
  \item Without any gap, $\|\theta(T)-\theta_0\|\le TM\|r_0\|/m$.
\end{enumerate}
\end{theorem}

\begin{proof}
  \uses{thm:ntkgf:displacementLinear, thm:ntkgf:expDecayIcc, lem:ntkf:gronwallBound}
  \leanok
By Theorem~\ref{thm:ntkgf:mseGradient}, $\|\theta'(t)\|=\|\nabla L(\theta(t))\|\le\frac1m\|J\|\|r(t)\|\le\frac Mm\|r(t)\|$, and Theorem~\ref{thm:ntkgf:expDecayIcc} gives $\|r(t)\|\le\|r_0\|e^{-\lambda t/m}$.
Theorem~\ref{thm:ntkgf:displacementLinear} with $\mathrm{Lin}=\mathrm{id}$ proves (1). For (2), $\int_0^Te^{-\lambda t/m}dt\le m/\lambda$ (Lemma~\ref{lem:ntkf:gronwallBound}), so the total distance traveled converges.
For (3), positive semidefiniteness alone makes $\|r(t)\|$ nonincreasing, so the speed is at most $\frac Mm\|r_0\|$, and the bound grows linearly in $T$ as it must.
\end{proof}

\begin{theorem}[Rayleigh bound on a ball]\label{thm:ntkgf:rayleighBall}
  \uses{thm:ntkgf:rayleighStability}
  \lean{NTK.rayleigh_lower_bound_on_ball}
  \leanok
Suppose on the closed ball $\{\|\theta-\theta_0\|\le r\}$ the Jacobian satisfies $\|J(\theta)\|\le M$ and $\|J(\theta)-J(\theta_0)\|\le L_J\|\theta-\theta_0\|$, the initial kernel satisfies
the Rayleigh--Ritz condition with parameter $\lambda_0$, and $2ML_Jr\le\lambda_0/2$. Then every $\theta$ in the ball satisfies the Rayleigh--Ritz condition with parameter $\lambda_0/2$.
\end{theorem}

\begin{proof}
  \leanok
By Theorem~\ref{thm:ntkgf:rayleighStability}(4) the parameter is at least $\lambda_0-2ML_J\|\theta-\theta_0\|\ge\lambda_0-2ML_Jr\ge\lambda_0/2$.
\end{proof}

\begin{theorem}[Lazy-training displacement bound (the bootstrap)]\label{thm:ntkgf:lazyDisplacement}
  \uses{thm:ntkf:bootstrap, thm:ntkgf:rayleighBall, thm:ntkgf:displacementBound}
  \lean{NTK.lazy_training_displacement_bound}
  \leanok
Let $\theta$ be a forward gradient flow of the MSE loss from $\theta_0$ with differentiable outputs. Let $M,L_J\ge0$, $\lambda_0>0$, $0\le r$ and $C<r$ satisfy
\begin{itemize}
  \item $2ML_Jr\le\lambda_0/2$,\qquad $M\|r_0\|/(\lambda_0/2)\le C$,
  \item the initial kernel $K(\theta_0)$ satisfies the Rayleigh--Ritz condition with parameter $\lambda_0$,
  \item on the ball $\{\|\theta-\theta_0\|\le r\}$: $\|J(\theta)\|\le M$ and $\|J(\theta)-J(\theta_0)\|\le L_J\|\theta-\theta_0\|$.
\end{itemize}
Then $\|\theta(T)-\theta_0\|\le C$ for every $T\ge0$.
\end{theorem}

\begin{proof}
  \uses{thm:ntkf:bootstrap, thm:ntkgf:rayleighBall, thm:ntkgf:displacementBound}
  \leanok
Apply the bootstrap principle (Theorem~\ref{thm:ntkf:bootstrap}) to $d(t)=\|\theta(t)-\theta_0\|$, which is continuous with $d(0)=0\le r$. The core claim is: if $d\le r$ on $[0,S]$, then
$d(S)\le C$. Indeed, inside the ball Theorem~\ref{thm:ntkgf:rayleighBall} gives the Rayleigh--Ritz parameter $\lambda_0/2$ along $[0,S]$ and the Jacobian is $M$-bounded, so
Theorem~\ref{thm:ntkgf:displacementBound}(2) gives $d(S)\le M\|r_0\|/(\lambda_0/2)\le C$. Since $C<r$, the bootstrap concludes $d\le C$ for all times, in particular the trajectory
never leaves the ball.
\end{proof}

\begin{theorem}[Global consequences of the ball hypotheses]\label{thm:ntkgf:lazyGlobal}
  \uses{thm:ntkgf:lazyDisplacement, thm:ntkgf:expDecayIcc, thm:ntkgf:kernelLipschitz}
  \lean{NTK.lazy_training_global_bounds_of_ball_hypotheses,
        NTK.lazy_training_kernel_freeze_bound_of_ball_hypotheses}
  \leanok
Under the hypotheses of Theorem~\ref{thm:ntkgf:lazyDisplacement}, for every $t\ge0$:
(i) $\|\theta(t)-\theta_0\|\le C$;
(ii) $K(\theta(t))$ satisfies the Rayleigh--Ritz condition with parameter $\lambda_0/2$;
(iii) $\|K(\theta(t))-K(\theta_0)\|\le(2ML_J)\,C$;
(iv) $\|r(t)\|\le\|r_0\|\,e^{-\lambda_0t/(2m)}$;
(v) $L(\theta(t))\le L(\theta_0)\,e^{-\lambda_0t/m}$.
In particular no free hypotheses on the displacement or the kernel Lipschitz bound remain: both are derived.
\end{theorem}

\begin{proof}
  \uses{thm:ntkgf:lazyDisplacement, thm:ntkgf:expDecayIcc, thm:ntkgf:kernelLipschitz}
  \leanok
(i) is Theorem~\ref{thm:ntkgf:lazyDisplacement}; (ii) follows from (i) and Theorem~\ref{thm:ntkgf:rayleighBall}; (iii) from (i) and Theorem~\ref{thm:ntkgf:kernelLipschitz}. For (iv) and (v)
the residual satisfies $\partial_tr=-\frac1mK(\theta(t))r$ (Theorem~\ref{thm:ntkgf:outputODE}) and the kernel satisfies the Rayleigh--Ritz condition with parameter $\lambda_0/2$ on $[0,t]$ by (ii), so
Theorem~\ref{thm:ntkgf:expDecayIcc} applies on $[0,t]$ for every $t$.
\end{proof}

\begin{theorem}[Finite horizon without a spectral gap]\label{thm:ntkgf:finiteHorizon}
  \uses{thm:ntkf:bootstrap, thm:ntkgf:displacementBound, thm:ntkgf:kernelLipschitz}
  \lean{NTK.finite_horizon_displacement_bound, NTK.finite_horizon_kernel_freeze_bound}
  \leanok
Let $\theta$ be a forward gradient flow of the MSE loss from $\theta_0$, $T,M\ge0$, $0\le r$, and $C<r$ with $TM\|r_0\|/m\le C$. If $\|J(\theta)\|\le M$ on the ball of radius $r$ about $\theta_0$, then
$\|\theta(t)-\theta_0\|\le C$ for all $t\in[0,T]$. If in addition $\|J(\theta)-J(\theta_0)\|\le L_J\|\theta-\theta_0\|$ on the ball, then $\|K(\theta(t))-K(\theta_0)\|\le(2ML_J)C$ for $t\in[0,T]$.
\end{theorem}

\begin{proof}
  \uses{thm:ntkf:bootstrap, thm:ntkgf:displacementBound, thm:ntkgf:kernelLipschitz}
  \leanok
Same bootstrap as Theorem~\ref{thm:ntkgf:lazyDisplacement}, with the core claim supplied by the no-gap bound Theorem~\ref{thm:ntkgf:displacementBound}(3): inside the ball on $[0,S]$,
$d(S)\le SM\|r_0\|/m\le C$. The kernel statement is Theorem~\ref{thm:ntkgf:kernelLipschitz}.
\end{proof}

\begin{remark}[Role of this section in the NTK development]\label{rmk:ntkgf:lazyRole}
  \uses{thm:ntkgf:displacementBound, thm:ntkgf:lazyDisplacement, thm:ntkgf:lazyGlobal, thm:ntkgf:finiteHorizon, thm:ntkf:bootstrap}
\emph{Main results.} The section answers one question: when does gradient flow keep the parameters, and hence the kernel, near their initial values? There are two answers, depending on whether the initial kernel has a spectral gap.
\begin{itemize}
  \item \emph{With a gap} (Theorems~\ref{thm:ntkgf:lazyDisplacement} and~\ref{thm:ntkgf:lazyGlobal}). If $K(\theta_0)$ has smallest eigenvalue at least $\lambda_0>0$ and the Jacobian is bounded and Lipschitz on a ball around $\theta_0$ with $2ML_Jr\le\lambda_0/2$, then the flow never leaves the ball: $\|\theta(t)-\theta_0\|\le C$ for all $t\ge0$, for any $C<r$ with $C\ge2M\|r_0\|/\lambda_0$. As consequences the kernel keeps the gap $\lambda_0/2$, moves by at most $2ML_JC$, and the residual and loss decay exponentially, all uniformly in time.
  \item \emph{Without a gap} (Theorem~\ref{thm:ntkgf:finiteHorizon}). Positive semidefiniteness alone only makes $\|r(t)\|$ nonincreasing, so the displacement grows linearly, $\|\theta(t)-\theta_0\|\le tM\|r_0\|/m$. The same bootstrap then gives the same conclusions on a fixed horizon $[0,T]$, as long as the ball radius is chosen larger than $C=TM\|r_0\|/m$. The price is that nothing is uniform in time and no decay of the residual is obtained.
\end{itemize}

\emph{What it builds on.}
\begin{itemize}
  \item The kernel-stability section supplies the two analytic inputs: Lipschitz propagation from $J$ to $K$ (Theorem~\ref{thm:ntkgf:kernelLipschitz}) and the persistence of a Rayleigh--Ritz gap under small displacement (Theorem~\ref{thm:ntkgf:rayleighStability}, packaged on a ball in Theorem~\ref{thm:ntkgf:rayleighBall}).
  \item The dynamics sections supply the speed of the flow. The gradient formula (Theorem~\ref{thm:ntkgf:mseGradient}) gives $\|\theta'\|\le\frac Mm\|r\|$, and the residual bounds of Theorems~\ref{thm:ntkgf:expDecay} and~\ref{thm:ntkgf:expDecayIcc} (with a gap) or Lemma~\ref{lem:ntkgf:lossMonotone} (without) control $\|r\|$. Theorem~\ref{thm:ntkgf:displacementBound} integrates the two, and it is the only place where the two cases differ.
  \item The bootstrap principle (Theorem~\ref{thm:ntkf:bootstrap}) closes the circularity: the estimates hold only inside the ball, and the ball is exactly what has to be proved. Both main theorems use it in the same way, and only the core estimate $d(S)\le C$ changes.
\end{itemize}

\emph{What it is used for.} This section is deterministic and takes the Jacobian estimates as hypotheses. It does not say that they hold for a wide network.
\begin{itemize}
  \item The next section turns the drift bound $2ML_JC$ into a limit as the width grows (Theorem~\ref{thm:ntkgf:kernelFreeze}), once the Jacobian estimates carry the appropriate powers of $n$.
  \item In the two-layer chapter the Jacobian estimates are proved with high probability over the initialization (Theorem~\ref{thm:ntktl:ballBounds}). Combining them with Theorem~\ref{thm:ntkgf:lazyGlobal} gives the uniform-in-time good event and the freezing of the kernel for all $t\ge0$ in the spectral-gap case. Combining them with Theorem~\ref{thm:ntkgf:finiteHorizon} gives the corresponding finite-horizon event (Theorem~\ref{thm:ntktl:finiteHorizonEvent}) and the finite-horizon stationarity results (Theorem~\ref{thm:ntktl:finiteHorizonLimits}), which need no gap and so also cover cases where the initial kernel is degenerate.
  \item The block-speed bound of Theorem~\ref{thm:ntkgf:displacementLinear}(2) is the ingredient for the per-neuron displacement estimate (Theorem~\ref{thm:ntktl:neuronDisplacement}).
  \item The kernel drift obtained here is the quantity $\varepsilon_K$ that enters the frozen-kernel comparison of Theorem~\ref{thm:ntkgf:residualFrozen}, which is how ``the kernel barely moves'' becomes ``the trained network follows the closed-form linear dynamics''.
\end{itemize}
\end{remark}


\section{Infinite-width asymptotics}\label{sec:ntkgf:infiniteWidth}

\begin{lemma}[Elementary limits]\label{lem:ntkgf:limits}
  \lean{NTK.tendsto_const_div_sqrt_nat_atTop, NTK.tendsto_chebyshev_bound_atTop,
        NTK.tendsto_lazy_training_kernel_freeze, NTK.tendsto_shallowEmpiricalNTK_chebyshev_bound}
  \leanok
For constants $c,C,\varepsilon$: $c/\sqrt n\to0$ and $C/(\varepsilon^2n)\to0$ as $n\to\infty$; consequently $L_KC/\sqrt n\to0$ and $\mathrm{ofReal}(C/(\varepsilon^2n))\to0$ in $[0,\infty]$.
\end{lemma}

\begin{theorem}[Kernel constancy (lazy training)]\label{thm:ntkgf:kernelFreeze}
  \uses{thm:ntkgf:kernelLipschitz, lem:ntkgf:limits}
  \lean{NTK.lazy_training_kernel_freeze_bound,
        NTK.empiricalNTKMatrix_trajectory_freeze_of_jacobian_bound,
        NTK.tendsto_lazy_training_kernel_freeze_matrix}
  \leanok
\begin{enumerate}
  \item If $\|\theta(t)-\theta_0\|\le C/\sqrt n$ for all $t\ge0$ and the empirical NTK matrix is $L_K$-Lipschitz along the trajectory, then
        $\|K(\theta(t))-K(\theta_0)\|\le L_KC/\sqrt n$ for all $t\ge0$.
  \item If $\|\theta(t)-\theta_0\|\le C$, $\|J(\theta(t))\|\le M$, $\|J(\theta_0)\|\le M$ and $J$ is $L_J$-Lipschitz along the trajectory, then
        $\|K(\theta(t))-K(\theta_0)\|\le2ML_JC$. (The width decay $n^{-1/2}$ belongs on $L_J$, coming from the $n^{-1/2}$ in the parametrization, not on the displacement.)
  \item For a family of models $f^{(n)}$ with $\|\theta^{(n)}(t)-\theta^{(n)}_0\|\le C/\sqrt n$ and $L_K$-Lipschitz kernel,
        $\|K^{(n)}(\theta^{(n)}(t))-K^{(n)}(\theta^{(n)}_0)\|\to0$ as $n\to\infty$, for each fixed $t\ge0$.
\end{enumerate}
\end{theorem}

\begin{proof}
  \uses{thm:ntkgf:kernelLipschitz, thm:ntkgf:displacementBound, thm:ntkgf:lazyDisplacement, lem:ntkgf:limits}
  \leanok
(1) is the Lipschitz bound composed with the displacement bound: the Lipschitz property of the kernel in the parameters (the kind of estimate given by Theorem~\ref{thm:ntkgf:kernelLipschitz}) turns $\|\theta(t)-\theta_0\|\le C/\sqrt n$ into $\|K(\theta(t))-K(\theta_0)\|\le L_K\|\theta(t)-\theta_0\|\le L_KC/\sqrt n$, and the displacement bound is the hypothesis that Theorems~\ref{thm:ntkgf:displacementBound} and~\ref{thm:ntkgf:lazyDisplacement} establish for the gradient flow. (2) is Theorem~\ref{thm:ntkgf:kernelLipschitz}; (3) follows from (1) and Lemma~\ref{lem:ntkgf:limits}.
\end{proof}

\begin{theorem}[Deterministic initialization]\label{thm:ntkgf:deterministicInit}
  \uses{def:ntk:shallowEmpiricalNTK, lem:ntkhelp:symmetryPsd, lem:ntkhelp:frobGradient, def:ntk:shallowLimitingNTK, lem:ntk:ntkConverges, thm:ntkf:slln, lem:ntkf:chebyshev}
  \lean{NTK.deterministic_initialization_shallowEmpiricalNTK_tendsto_ae,
        NTK.deterministic_initialization_empiricalNTKMatrix_tendsto_ae,
        NTK.deterministic_initialization_chebyshev_bound}
  \leanok
Let $\sigma'$ be measurable and bounded, and rows $w_0,w_1,\dots$ i.i.d.\ $\mathcal N(0,I_d)$. Almost surely the empirical kernel $k_n(x,x')$ built from the first $n$ rows (Definition~\ref{def:ntk:shallowEmpiricalNTK}) converges to the limiting NTK $k_\infty(x,x')$ (Definition~\ref{def:ntk:shallowLimitingNTK}) entrywise, and for a finite dataset the empirical Gram matrix $\bigl(k_n(x^\alpha,x^\beta)\bigr)_{\alpha\beta}$ (the matrix of Lemma~\ref{lem:ntkhelp:symmetryPsd}(2), a Gram matrix of gradients by Lemma~\ref{lem:ntkhelp:frobGradient}) converges entrywise to $\bigl(k_\infty(x^\alpha,x^\beta)\bigr)_{\alpha\beta}$.
Quantitatively, if $k_n(x,x')\in L^2$ and $\operatorname{Var}(k_n(x,x'))\le C/n$, then for every $\varepsilon>0$
\[
  \Pr\Bigl[\,\bigl|k_n(x,x')-\mathbb E_{w_0,\dots,w_{n-1}}\bigl[k_n(x,x')\bigr]\bigr|\ge\varepsilon\,\Bigr]\le\frac{C}{\varepsilon^2n},
\]
where $\mathbb E_{w_0,\dots,w_{n-1}}$ is the expectation over the $n$ i.i.d.\ Gaussian rows $w_0,\dots,w_{n-1}\sim\mathcal N(0,I_d)$ that define $k_n$ (the inputs $x,x'$ are fixed), and $\Pr$ is the probability over the same rows.
\end{theorem}

\begin{proof}
  \uses{lem:ntk:ntkConverges, lem:ntkf:chebyshev}
  \leanok
The almost sure statements are Lemma~\ref{lem:ntk:ntkConverges} applied entrywise (finitely many entries, so ``almost surely for each'' implies ``almost surely for all''). The tail bound is Chebyshev's inequality
(Lemma~\ref{lem:ntkf:chebyshev}) with the assumed variance bound.
\end{proof}

\begin{remark}[Role of this section in the NTK development]\label{rmk:ntkgf:infiniteWidthRole}
  \uses{thm:ntkgf:kernelFreeze, thm:ntkgf:deterministicInit, lem:ntkgf:limits, thm:ntkgf:kernelLipschitz}
\emph{Main results.} The previous section is deterministic: for one network it bounds how far the kernel moves in terms of the displacement of the parameters. This section adds the width $n$ and shows the two halves of the infinite-width limit.
\begin{itemize}
  \item \emph{The kernel does not move} (Theorem~\ref{thm:ntkgf:kernelFreeze}). If the displacement is $O(n^{-1/2})$ and the kernel is Lipschitz in the parameters, then $\|K(\theta(t))-K(\theta_0)\|\le L_KC/\sqrt n\to0$ for every $t\ge0$. Part~(2) is the version used when the width decay sits on the Jacobian Lipschitz constant $L_J$ rather than on the displacement.
  \item \emph{The initial kernel converges} (Theorem~\ref{thm:ntkgf:deterministicInit}). At random initialization the empirical kernel $k_n$ converges almost surely to the deterministic limiting NTK $k_\infty$, with a Chebyshev tail bound of order $C/(\varepsilon^2n)$.
\end{itemize}
Together they say that a wide network trains with a kernel that starts near $K_\infty$ and stays there, which is exactly the frozen-kernel hypothesis of the closed-form dynamics.

\emph{What it builds on.}
\begin{itemize}
  \item Kernel constancy is the Lipschitz propagation of Theorem~\ref{thm:ntkgf:kernelLipschitz} combined with a displacement bound from the lazy-training section (Theorems~\ref{thm:ntkgf:displacementBound} and~\ref{thm:ntkgf:lazyDisplacement}). Lemma~\ref{lem:ntkgf:limits} supplies the elementary limits $c/\sqrt n\to0$ and $C/(\varepsilon^2n)\to0$.
  \item Deterministic initialization is the strong law of large numbers and Chebyshev's inequality (Theorem~\ref{thm:ntkf:slln}, Lemma~\ref{lem:ntkf:chebyshev}) applied to the empirical NTK of the shallow network (Definition~\ref{def:ntk:shallowEmpiricalNTK}, Lemma~\ref{lem:ntk:ntkConverges}). The dataset-level Gram matrix is the one of Lemma~\ref{lem:ntkhelp:symmetryPsd}(2).
  \item Both statements concern the objects of the gradient-flow chapter, but neither uses the dynamics: kernel constancy takes the displacement bound as a hypothesis, and initialization only concerns the law of $\theta_0$.
\end{itemize}

\emph{What it is used for.} The results here are stated for a generic family of models. They are the template that the two-layer chapter instantiates with probabilities and explicit constants.
\begin{itemize}
  \item The convergence of the initial kernel is made quantitative for the two-layer network in Theorems~\ref{thm:ntktl:ntkLLN} and~\ref{thm:ntktl:ntkChebyshev}. The Chebyshev bound is what turns $K_n(0)\to K_\infty$ into an initial spectral gap with high probability (Theorem~\ref{thm:ntktl:spectralGapInit}), the input needed by the gap case of the lazy-training section.
  \item Kernel constancy is proved for Gaussian initialization with all hypotheses discharged in Theorem~\ref{thm:ntktl:kernelFreezeGaussian}, and the rate $O(n^{-1/2})$ is obtained from a moment bound on a single neuron in Theorems~\ref{thm:ntktl:driftRate} and~\ref{thm:ntktl:globalEventInvSqrt}.
  \item Both limits feed the comparison with the frozen-kernel dynamics: $K_n(0)\to K_\infty$ fixes the reference curve $\exp(-\frac tmK_\infty)r_0$, and kernel constancy supplies the small $\varepsilon_K$ in Theorem~\ref{thm:ntkgf:residualFrozen}. The next sections solve the frozen dynamics in the fixed-feature setting (Theorem~\ref{thm:ntkgf:affineFlow}) and compare test-point predictions (Theorem~\ref{thm:ntkgf:predictionError}).
\end{itemize}
\end{remark}


\section{Fixed-feature (affine) dynamics and the minimum-norm interpolant}\label{sec:ntkgf:affine}

Consider the affine model $\theta\mapsto f_0+Jw$ with fixed Jacobian $J\in\mathbb R^{m\times P}$ and residual $r_0+Jw$ (here $r_0=f_0-y$), trained by gradient flow from $w(0)=0$:
\begin{equation}\label{eq:ntkgf:affineFlow}
  w'(t)=-\frac1m\,J^\top\bigl(r_0+Jw(t)\bigr).
\end{equation}

\begin{theorem}[Closed form of the affine flow]\label{thm:ntkgf:affineFlow}
  \uses{thm:ntkf:uniqueness, thm:ntkgf:matrixExp}
  \lean{NTK.hasDerivAt_affineFlowSolution, NTK.affineFlow_eq_solution}
  \leanok
Assume the Gram matrix $JJ^\top$ is invertible. The curve
\[
  w(t)=-J^\top(JJ^\top)^{-1}\Bigl(r_0-\exp\!\bigl(-\tfrac tmJJ^\top\bigr)r_0\Bigr)
\]
solves~\eqref{eq:ntkgf:affineFlow} with $w(0)=0$. Any forward solution of~\eqref{eq:ntkgf:affineFlow} continuous on $[0,\infty)$ with $w(0)=0$ equals this curve for $t\ge0$; in particular $w(t)$ lies in the range of $J^\top$.
\end{theorem}

\begin{proof}
  \uses{thm:ntkf:uniqueness, thm:ntkgf:matrixExp}
  \leanok
Differentiate: $w'(t)=-J^\top(JJ^\top)^{-1}\cdot\frac1mJJ^\top\exp(-\frac tmJJ^\top)r_0=-\frac1mJ^\top\exp(-\frac tmJJ^\top)r_0$, while $r_0+Jw(t)=\exp(-\frac tmJJ^\top)r_0$ since $J w(t)=-(r_0-\exp(\cdots)r_0)$ ---
so both sides of~\eqref{eq:ntkgf:affineFlow} agree. Uniqueness is Theorem~\ref{thm:ntkf:uniqueness}, because the vector field $w\mapsto-\frac1mJ^\top(r_0+Jw)$ is affine, hence Lipschitz.
\end{proof}

\begin{theorem}[Minimum-norm characterization]\label{thm:ntkgf:minNorm}
  \uses{lem:ntkf:matrixCLM}
  \lean{NTK.minNorm_pythagoras, NTK.affine_minNorm_pythagoras, NTK.norm_affine_limit_le,
        NTK.eq_affine_limit_of_norm_le}
  \leanok
Let $A:E\to F$ be a bounded linear map between real Hilbert spaces and suppose $w_\infty=A^\dagger a$ interpolates, $Aw_\infty=b$. Then every $w'$ with $Aw'=b$ satisfies the Pythagoras identity
$\|w'\|^2=\|w_\infty\|^2+\|w'-w_\infty\|^2$; hence $w_\infty$ is the unique minimum-norm interpolant.
In particular, for $J$ with invertible Gram matrix and $w_\infty=-J^\top(JJ^\top)^{-1}r_0$: $Jw_\infty=-r_0$, every $w'$ with $Jw'=-r_0$ satisfies
$\|w'\|^2=\|w_\infty\|^2+\|w'-w_\infty\|^2$, so $\|w_\infty\|\le\|w'\|$ and equality of norms forces $w'=w_\infty$.
\end{theorem}

\begin{proof}
  \uses{lem:ntkf:matrixCLM}
  \leanok
$\|w'\|^2=\|w_\infty+(w'-w_\infty)\|^2=\|w_\infty\|^2+2\langle w_\infty,w'-w_\infty\rangle+\|w'-w_\infty\|^2$, and $\langle A^\dagger a,w'-w_\infty\rangle=\langle a,A(w'-w_\infty)\rangle=0$ since both interpolate. For the matrix case take
$A=\operatorname{matrixCLM}(J)$, whose adjoint is $\operatorname{matrixCLM}(J^\top)$ (Lemma~\ref{lem:ntkf:matrixCLM}), and $a=-(JJ^\top)^{-1}r_0$.
\end{proof}

\begin{theorem}[Convergence to the interpolant and kernel regression]\label{thm:ntkgf:affineLimit}
  \uses{thm:ntkgf:affineFlow, thm:ntkgf:minNorm, thm:ntkf:spectralGap, def:ntkgf:crossKernel}
  \lean{NTK.tendsto_affineFlowSolution, NTK.inner_affine_limit}
  \leanok
If $JJ^\top$ is positive definite, the affine flow converges, $w(t)\to w_\infty=-J^\top(JJ^\top)^{-1}r_0$, at the exponential rate $\lambda_{\min}(JJ^\top)/m$; the limit interpolates and has least norm. For any feature vector $g$,
\[
  \langle g,w_\infty\rangle=-(Jg)^\top(JJ^\top)^{-1}r_0 .
\]
With $J=J(\theta_0)$, $g=\nabla f(x;\theta_0)$ and $r_0=f_0(X)-y$ this equals $k_0(x,X)^\top K_0^{-1}\bigl(y-f_0(X)\bigr)$, where $k_0(x,X)=J(\theta_0)g=k_{\theta_0}(x)$ is the initial train--test cross-kernel (Definition~\ref{def:ntkgf:crossKernel}, Lemma~\ref{lem:ntkgf:crossKernelFormulas}) and $K_0=K(\theta_0)$: the frozen-feature minimum-norm predictor is the kernel-regression formula with the empirical kernels.
\end{theorem}

\begin{proof}
  \uses{thm:ntkgf:affineFlow, thm:ntkgf:minNorm, thm:ntkf:spectralGap}
  \leanok
By Theorem~\ref{thm:ntkf:spectralGap}, $JJ^\top\succeq cI$ for some $c>0$, so $\exp(-\frac tmJJ^\top)r_0\to0$ at rate $c/m$ and the closed form of Theorem~\ref{thm:ntkgf:affineFlow} converges to $-J^\top(JJ^\top)^{-1}r_0$.
For the last identity, $\langle g,J^\top v\rangle=\langle Jg,v\rangle$.
\end{proof}


\section{Test-point prediction along a gradient flow}\label{sec:ntkgf:testPrediction}

This section asks how the trained network behaves at a \emph{test} input $x$. The ingredients are:
\begin{itemize}
  \item the \emph{frozen tangent feature} $g:=\nabla_\theta f(x;\theta_0)$ of the test input at the initialization, and the initial residual $r_0=r(\theta_0)$;
  \item the empirical cross-kernel at initialization, $k_{\theta_0}(x)=J(\theta_0)\,g$ (Definition~\ref{def:ntkgf:crossKernel} and Lemma~\ref{lem:ntkgf:crossKernelFormulas}), and the cross-kernel $k_\theta(x)=J(\theta)\,\nabla f(x;\theta)$ along the flow, which is what actually drives the test output (Remark~\ref{rmk:ntkgf:crossKernelMeaning});
  \item a \emph{limiting} kernel matrix $K_\infty\in\mathbb R^{m\times m}$ (symmetric positive semidefinite) and a \emph{limiting cross-kernel} vector $k_\infty\in\mathbb R^m$, the deterministic counterpart of the empirical cross-kernel $k_\theta(x)$ of Definition~\ref{def:ntkgf:crossKernel}. In the infinite-width limit these are the NTK of the training points and the NTK between the test input and the training points, respectively; here they are just given deterministic objects, and the hypotheses below say how close the empirical quantities are to them;
  \item a vector $a$ with $K_\infty a=k_\infty$, the \emph{kernel-regression weights} ($a=K_\infty^{-1}k_\infty$ when $K_\infty$ is invertible; no invertibility is required here).
\end{itemize}
If the kernels were exactly frozen at $K_\infty$ and $k_\infty$, the linearized test output would move by $-a^\top\bigl(r_0-\exp(-\frac tmK_\infty)r_0\bigr)$ (compare Theorem~\ref{thm:ntkgf:affineLimit}, whose limit is the kernel-regression formula, and Theorem~\ref{thm:ntkgf:matrixExp} for the residual $\exp(-\frac tmK_\infty)r_0$). The \emph{prediction error} at time $t$ compares the actual displacement of the parameter in the test direction $g$ with this frozen-kernel prediction:
\[
  e(t):=\langle g,\theta(t)-\theta_0\rangle+a^\top\Bigl(r_0-\exp\!\bigl(-\tfrac tmK_\infty\bigr)r_0\Bigr).
\]
It starts at $e(0)=0$. The deterministic estimates below bound $|e(t)|$ in terms of three quantities that the probabilistic chapters then show to be small: the Jacobian drift $\|J(\theta(s))-J(\theta_0)\|\le\varepsilon_J$, the kernel drift $\|K(\theta(s))-K_\infty\|\le\varepsilon_K$, and the initial cross-kernel error $\|J(\theta_0)g-k_\infty\|=\|k_{\theta_0}(x)-k_\infty\|\le\varepsilon_k$ (Definition~\ref{def:ntkgf:crossKernel}, Lemma~\ref{lem:ntkgf:crossKernelFormulas}(2)).
For the two-layer network, convergence of the initial cross-kernel (Definition~\ref{def:ntkgf:crossKernel}) to $k_\infty$ is Theorem~\ref{thm:ntktl:crossKernelInit}, its freezing along the flow is Theorem~\ref{thm:ntktl:crossKernelFreeze}, and the resulting test-prediction limit is Theorem~\ref{thm:ntktl:testPrediction}.

\begin{lemma}[Derivative of the prediction error]\label{lem:ntkgf:predictionDerivative}
  \uses{thm:ntkgf:outputODE, thm:ntkgf:eigenmodes, def:ntkgf:crossKernel, lem:ntkgf:crossKernelFormulas}
  \lean{NTK.hasDerivAt_inner_tangentFeature_displacement, NTK.hasDerivAt_frozenPrediction,
        NTK.abs_dotProduct_le_norm_mul_norm, NTK.hasDerivAt_predictionError_abs_le}
  \leanok
Along a gradient flow, $\frac d{dt}\langle g,\theta(t)-\theta_0\rangle=-\frac1m\sum_\alpha r^\alpha\langle g,\nabla f(x^\alpha;\theta(t))\rangle$ and
$\frac{d}{dt}\bigl[-a^\top(r_0-\exp(-\tfrac tmK_\infty)r_0)\bigr]=-\frac1mk_\infty^\top\exp(-\tfrac tmK_\infty)r_0$. Consequently $e$ is differentiable at every time $c$ at which the flow equation holds, with
\[
  |e'(c)|\le\frac1m\Bigl(\|r(c)\|\bigl(\|g\|\,\|J(\theta(c))-J(\theta_0)\|+\varepsilon_k\bigr)+\|k_\infty\|\,\bigl\|r(c)-\exp(-\tfrac cmK_\infty)r_0\bigr\|\Bigr),
\]
where $\varepsilon_k\ge\|J(\theta_0)g-k_\infty\|$. The three terms are the Jacobian drift, the cross-kernel error (how far the empirical initial cross-kernel $k_{\theta_0}(x)=J(\theta_0)g$ of Definition~\ref{def:ntkgf:crossKernel} is from the limiting cross-kernel $k_\infty$), and the residual-versus-frozen-exponential error.
\end{lemma}

\begin{proof}
  \leanok
$\frac{d}{dt}\langle g,\theta(t)\rangle=\langle g,\theta'\rangle=-\frac1m\langle g,J^\top r\rangle=-\frac1m(Jg)^\top r$ and, using $K_\infty a=k_\infty$, the frozen part has derivative $-\frac1m k_\infty^\top\exp(\cdots)r_0$.
Adding: $e'=-\frac1m\bigl[(J(\theta(c))g)^\top r(c)-k_\infty^\top\exp(-\tfrac cmK_\infty)r_0\bigr]$. Write
$(J(\theta(c))g)^\top r(c)-k_\infty^\top s(c)=\bigl((J(\theta(c))-J(\theta_0))g\bigr)^\top r(c)+(J(\theta_0)g-k_\infty)^\top r(c)+k_\infty^\top(r(c)-s(c))$ with $s(c)=\exp(-\tfrac cmK_\infty)r_0$ and apply Cauchy--Schwarz to each term.
\end{proof}

\begin{theorem}[Deterministic test-point prediction error]\label{thm:ntkgf:predictionError}
  \uses{lem:ntkgf:predictionDerivative, thm:ntkgf:residualFrozen, def:ntkgf:crossKernel, lem:ntkgf:crossKernelFormulas}
  \lean{NTK.abs_inner_displacement_add_frozenPrediction_le}
  \leanok
Let $\theta$ be a forward gradient flow of the MSE loss from $\theta_0$ and suppose on $[0,T]$: $\|J(\theta(s))-J(\theta_0)\|\le\varepsilon_J$ and $\|K(\theta(s))-K_\infty\|\le\varepsilon_K$, and also
$\|J(\theta_0)g-k_\infty\|\le\varepsilon_k$, i.e.\ the empirical initial cross-kernel $k_{\theta_0}(x)=J(\theta_0)g$ (Definition~\ref{def:ntkgf:crossKernel}) is within $\varepsilon_k$ of the limiting cross-kernel $k_\infty$. Then for $t\in[0,T]$,
\[
  |e(t)|\le t\cdot\frac1m\Bigl(\|g\|\,\varepsilon_J\,\|r_0\|+\varepsilon_k\,\|r_0\|+\|k_\infty\|\cdot\tfrac1m\varepsilon_K\|r_0\|\,T\Bigr).
\]
\end{theorem}

\begin{proof}
  \uses{lem:ntkgf:predictionDerivative, thm:ntkgf:residualFrozen, def:ntkgf:crossKernel}
  \leanok
Since $e(0)=0$, it suffices to integrate the bound of Lemma~\ref{lem:ntkgf:predictionDerivative}. On $[0,T]$ one has $\|r(s)\|\le\|r_0\|$ (the loss is nonincreasing), and Theorem~\ref{thm:ntkgf:residualFrozen}(2) gives
$\|r(s)-\exp(-\tfrac smK_\infty)r_0\|\le\frac1m\varepsilon_K\|r_0\|s\le\frac1m\varepsilon_K\|r_0\|T$.
\end{proof}

\begin{lemma}[Exponentially decaying derivative gives a small tail]\label{lem:ntkgf:expTail}
  \lean{NTK.abs_sub_le_of_abs_deriv_le_exp}
  \leanok
If $e$ is continuous on $[S,\infty)$, differentiable on $(S,\infty)$ with $|e'(s)|\le B e^{-\nu s}$, $\nu>0$, then for every $t\ge S$, $|e(t)-e(S)|\le\frac B\nu e^{-\nu S}$.
\end{lemma}

\begin{proof}
  \leanok
$|e(t)-e(S)|\le\int_S^t|e'|\le B\int_S^\infty e^{-\nu s}ds=\frac B\nu e^{-\nu S}$.
\end{proof}

\begin{theorem}[Uniform-in-time prediction error]\label{thm:ntkgf:predictionUniform}
  \uses{thm:ntkgf:predictionError, lem:ntkgf:expTail, def:ntkgf:crossKernel}
  \lean{NTK.abs_inner_displacement_add_frozenPrediction_le_of_exp_decay}
  \leanok
In addition to the hypotheses of Theorem~\ref{thm:ntkgf:predictionError} on a window $[0,S]$ (including the bound $\varepsilon_k$ on the initial cross-kernel error of Definition~\ref{def:ntkgf:crossKernel}), assume: the Jacobian drift is at most $\varepsilon_J$ for \emph{all} times; the residual decays as $\|r(s)\|\le\|r_0\|e^{-\nu s}$ with $\nu>0$;
and $K_\infty$ has spectral gap $\lambda$ with $\nu\le\lambda/m$. Then for every $t\ge0$,
\begin{multline*}
  |e(t)|\le S\cdot\frac1m\Bigl(\|g\|\varepsilon_J\|r_0\|+\varepsilon_k\|r_0\|+\|k_\infty\|\tfrac1m\varepsilon_K\|r_0\|S\Bigr)\\
  \qquad+\frac{\tfrac1m\|r_0\|\bigl(\|g\|\varepsilon_J+\varepsilon_k+2\|k_\infty\|\bigr)}{\nu}\,e^{-\nu S}.
\end{multline*}
\end{theorem}

\begin{proof}
  \uses{thm:ntkgf:predictionError, lem:ntkgf:expTail}
  \leanok
On $[0,S]$ use Theorem~\ref{thm:ntkgf:predictionError}. For $s>S$, both the actual residual and the frozen residual decay exponentially (the latter because $\nu\le\lambda/m$), so the three terms of
Lemma~\ref{lem:ntkgf:predictionDerivative} give $|e'(s)|\le\frac1m\|r_0\|(\|g\|\varepsilon_J+\varepsilon_k+2\|k_\infty\|)e^{-\nu s}$; Lemma~\ref{lem:ntkgf:expTail} bounds the tail.
\end{proof}
